% Draft data section, 2026-10-15 (Claude change 10-15, new file).
% Written to replace Section 4 of Flooding_Aug_2026.pdf after the ATTOM permit switch (10-02),
% the zone-aligned NFIP-ATTOM match (10-12) and the event-date / SFHA rules (10-15).
% Model: Wagner (2022, AEJ:Policy) Section III (one titled paragraph per administrative source,
% then sample restrictions, then the linkage with its limitation, then the summary table) and
% Gruber and Solomon (2026, NBER WP 35408), whose data paragraphs are short, say what each
% source is and what it is used for, and push institutional detail to the background section.
% Numbers marked [10-15] are from the 10-15 build (zone-aligned links, Oct 8 permit file, the
% complete.do / summary_table.do / elevations_by_state.do run of 10-15 15:41, after the cluster's
% permit status history landed; the morning build's numbers were replaced the same day).
% Script names, file paths and variable names are left out on purpose (data_revision_guide.md).

\section{Data}

We combine four administrative datasets: the universe of National Flood Insurance Program
(NFIP) policies and claims, FEMA's Hazard Mitigation Assistance (HMA) grant records, building
permits, and tax-assessor property records. We supplement these with FEMA's digital flood maps.
Each source contributes a distinct piece of the analysis: NFIP identifies the insured population
and records premiums, coverage and claims; HMA records federally funded elevations and their
cost; permits date privately and publicly funded elevations; and assessor records give property
values and a street address. The analysis covers single-family homes in Florida, Louisiana and
Texas from 2009 to 2025, deflated to 2023 dollars.

\paragraph{Flood insurance policies and claims.}
The NFIP policy data are the public OpenFEMA release of every policy transaction since 2009,
one record per policy-year, with premium, building and contents coverage, the rated flood zone,
construction year, pre- or post-FIRM status, primary-residence status and an elevated-building
indicator. The claims data record every paid claim with the date of loss and the building and
contents payouts; the multiple-loss file flags repetitive-loss and severe-repetitive-loss
structures. Addresses are redacted. Following \citet{wagner2022}, we identify a structure
approximately by its census block group, construction year and original policy date, which
yields a property-year panel. Where several structures share an identifier their claims are
pooled; we drop the 5,449 such identifiers (0.08 percent of properties) whose claims exceed
their coverage. The three states hold 7.1 million insured single-family structures, of which
3.3 million are rated in a Special Flood Hazard Area (SFHA) [10-15].

We use the NFIP elevated-building indicator only as a control. It behaves as a stock rather
than a flow: of 313,000 flagged structures nationally, about one percent ever change status
on the panel, and the elevation-certificate fields show that a change in the flag is a
certificate entering the rating, usually for a house that was already elevated, rather than a
construction event.\footnote{Within-structure comparisons of the certificate's lowest-floor
elevation before and after a flag change show a median rise of zero feet; only 9 percent rise
by three feet or more.} The flag therefore cannot date elevations, which is why we turn to
permits.

\paragraph{Mitigation grants.}
FEMA's HMA records cover every federally funded elevation project under the Flood Mitigation
Assistance (FMA), Hazard Mitigation Grant (HMGP), Pre-Disaster Mitigation, Building Resilient
Infrastructure and Communities, Repetitive Flood Claims and Severe Repetitive Loss programs.
The project file carries program, project type, approval status, benefit-cost ratio, number of
properties and federal obligation; the mitigated-properties file lists each treated property
with its ZIP code but no address, and includes applications that were never funded, which we
remove by requiring a funded project. Nationally the programs elevated about 18,800
single-family homes between 1997 and 2025, 70 percent of them under HMGP, at a median federal
cost of \$130,000 per home in 2023 dollars. Because properties carry no address, and grants FEMA never logged at the
property level carry no ZIP code, we measure grant exposure at the county level, as cumulative grant-elevated
homes and federal spending per insured home. Property-level grant lists are being obtained
from the states through records requests and will replace the county measure where available.

\paragraph{Building permits.}
Our universe of home elevations, funded or not, comes from ATTOM's building-permit product,
which compiles permits from local jurisdictions and attaches ATTOM's property identifier to
79 percent of them. We extract every permit in the three states (110 million) whose type,
subtype or description mentions an elevation, and screen the text for a house being raised,
removing elevation certificates, utility and equipment elevations, raised slabs on new
construction, boat lifts and demolitions. The screen leaves 17,143 addresses, of which 15,290
are retrofits of an existing house and the rest elevated new construction [10-15].\footnote{A
hand reading of the survivors in Duval County, Florida, found that the screen's remaining
errors are concentrated in one city's ``elevation certificate required'' stamp on unrelated
permits, which the final screen removes.} ATTOM's permit coverage is by jurisdiction and year:
within a reporting jurisdiction-year every permit is recorded, but jurisdictions enter the feed
at different dates. We therefore count all permits of any type by county, ZIP code and city by
year and treat a home's elevation status as observed only in county-years with a permit feed.
Florida is covered throughout our period; Texas from 2009 for about 90 percent of insured
homes; Louisiana from 2013 for 80 percent, with a gap in 2017--2018 in Jefferson Parish.

Each permit carries the jurisdiction's own status label and an issue date. ATTOM also
publishes a status history that resolves the labels into fifteen standard stages with effective
dates. We date an elevation at the year the permit was finaled; where the permit has a history
but never reached a completion stage, at the year of its last recorded stage; and otherwise at
the issue year. We keep an elevation only when the permit record shows it done: finaled,
completed or closed. Jurisdictions whose feeds never record a completion stage, which account
for a third of Louisiana's permits, are treated separately, since an ``issued'' label there is
the last label the feed ever writes.\footnote{The status history covers 86 to 96 percent of the elevation addresses in the three
states. Of the addresses, 51 percent in Florida, 38 in Louisiana and 47 in Texas have a completion
stage; about a quarter of those completion rows carry no date, so they count for completion but
not for the event year. The remaining addresses sit at issued, approved or applied, concentrated
in 2023 and later and in Louisiana feeds that never record completions. [10-15]}

\paragraph{Property records.}
Property characteristics come from ATTOM's tax-assessor file, which covers every parcel in the
three states with its street address, coordinates, land use, construction year, and assessed
and market values by tax year. We keep single-family residential parcels, geocode each to its
2010 and 2020 census block group, and join each parcel to its flood zone on FEMA's National
Flood Hazard Layer (NFHL). Market value is the assessor's market field, which we deflate to 2023
dollars and carry to each policy year from the nearest assessment.\footnote{Louisiana assesses
residential property at ten percent of market value, and in a handful of rural parishes the
market field sits at the assessed level; we leave those values in and note them in the
appendix.} Permits join to the assessor record on ATTOM's property identifier (93 to 97 percent
of linked permits) or, failing that, on the street address (85, 93 and 95 percent of elevation
addresses in Florida, Louisiana and Texas are linked) [10-15].

\subsection{Linking insured structures to properties}

NFIP's redaction means an insured structure cannot be matched to a parcel one to one. We
follow Wagner's (2022) cell match: an insured structure and a parcel are paired when they share
a census block group, construction year and flood zone, with the flood zone taken from the
policy's rating on the NFIP side and from the NFHL on the parcel side, placed on a common code
key.\footnote{NFIP rates zones AHB and AOB, A1--A30 and V1--V30 and the pre-1980 B and C
zones, which the current NFHL does not carry; we map each to its NFHL equivalent before
matching. Before this alignment, 15 percent of the panel could never reach a flood-zone
tier.} Where no parcel shares all three, the match relaxes in a fixed order: a five-year
construction band and FIRM status in place of the exact year, the ZIP code, the NFIP
community and the county in place of the block group, high- or low-risk in place of the zone,
and finally no zone at all; a last tier pairs structures with no construction year on block
group and zone alone. Within a cell the pairing is arbitrary, except that a parcel with an
elevation permit is paired first with a structure insured in the permit year.

The match assigns a parcel to 86 percent of insured structures in Florida, 65 percent in
Louisiana and 92 percent in Texas [10-15]. Louisiana's rate is low because 59 percent of its
single-family parcels have no construction year in the assessor file, so most of its matches
fall to the final tier. Among structures carrying an elevation permit, two thirds are matched
on the exact cell and 94 percent on a block-group tier of some kind [10-15]. Because the
pairing within a cell is arbitrary, property values and permits are in substance cell-level
measures: they describe the insured structure's immediate neighbourhood and vintage, not the
structure itself. The permit's location is exact; which of the cell's insured structures it is
attached to is not, and our event-study estimates should be read with that attenuation in mind.

\subsection{Sample}

The analysis sample is every NFIP-insured single-family structure in Florida, Louisiana and
Texas that is rated in an SFHA, observed in the policy years 2009 to 2025. We take SFHA status
from the policy's own rated zone and, for the few hundred structures per state whose policies
carry no zone, from the NFHL zone of the matched parcel. We restrict nothing else at the build
stage: the ATTOM match, the permit-feed coverage of the county-year and the status of the
permit ride along as flags, and each analysis applies the restriction it needs. The event
studies, for instance, use structures with an ATTOM match in counties with at least one
permitted home, and date treatment by the permit's event year. The sample holds 3,291,653
structures and 16.6 million property-years [10-15]. Three quarters of the structures are
matched to a parcel (77 percent in Florida, 56 in Louisiana, 87 in Texas), 92 percent sit in a
county-year with a permit feed (99, 69 and 92 percent), and 74 percent are both, which is the
set on which an elevation is observable. Among the 5,944 structures with a permitted retrofit,
43 percent have reached a completion stage in the status history (48 percent in Florida, 39 in
Louisiana and Texas); 470 of them have no issue date on the record and are dated by the history
alone [10-15].

Three features of the sample bear on interpretation. First, the panel starts in 2009, so claims
before that year are unobserved and a structure's flood history is left-censored. Second, a
structure enters the panel only while insured; elevations on homes that drop NFIP coverage are
unobserved after exit, and the cell match attaches them to another insured structure. Third,
permit-feed coverage is a property of the jurisdiction and year, so the elevation rate is
defined only on the covered subsample, which is nearly all of Florida and Texas but under
half of Louisiana's insured homes before 2013.

\subsection{Summary statistics}

Table~\ref{tab:summary} reports the sample [10-15]. The median structure pays \$644 a year
in premium and 1.2 percent of property-years carry a claim; claims paid exceed premiums by 7
percent over the period, 26 percent for pre-FIRM structures and 18 percent below premiums for
post-FIRM ones. Claims are skewed: among structures that ever claim, the median cumulative
payout is \$82,600, and \$104,000 among repetitive-loss structures. Fourteen percent of
structures are already rated as elevated when first observed, almost all of them post-FIRM
construction. Elevations during the period are rare: 5,944 permitted retrofits, or 1.1, 4.2
and 2.0 per 1,000 insured SFHA homes in Florida, Louisiana and Texas, an order of magnitude
more than the NFIP indicator shows changing, concentrated in the parishes and counties with the
largest repetitive-loss populations, at a median declared cost of \$132,000, about 56 percent
of the median property value. The
federal programs reach a similar number of homes: on the 09-20 build, grant-elevated homes
per insured SFHA home are 0.9 percent in Louisiana, 0.2 percent in Texas and 0.02 percent in
Florida.

\begin{table}[htbp]
\centering
\caption{Summary statistics, NFIP-insured single-family structures in SFHAs, FL LA TX, 2009--2025}
\label{tab:summary}
\begin{tabular}{lccc}
\hline
 & All & Pre-FIRM & Post-FIRM \\
\hline
\multicolumn{4}{l}{\textit{A. Insurance}} \\
\quad Premium (median, 2023 \$) & 644 & 888 & 574 \\
\quad Claims paid / premiums & 1.07 & 1.26 & 0.82 \\
\quad Share of property-years with a claim & 0.012 & 0.016 & 0.008 \\
\quad Cumulative claims if any (median, 2023 \$) & 82,573 & 96,483 & 45,856 \\
\multicolumn{4}{l}{\textit{B. Elevation}} \\
\quad Elevated at first observation (NFIP indicator) & 449,933 & 172,804 & 277,129 \\
\quad Permitted retrofit, 2009--2025 & 5,944 & 3,917 & 2,027 \\
\quad Declared retrofit cost (median, 2023 \$) & 132,001 & 138,181 & 110,787 \\
\multicolumn{4}{l}{\textit{C. Property}} \\
\quad Market value (median, 2023 \$) & 236,264 & 209,894 & 259,479 \\
\quad Construction year (median) & 1984 & 1970 & 1997 \\
\hline
Number of structures & 3,291,653 & 1,621,850 & 1,669,803 \\
Number of property-years & 16,590,573 & 7,462,964 & 9,127,609 \\
\hline
\end{tabular}
\begin{minipage}{0.95\linewidth}
\footnotesize \textit{Notes:} Premium, loss ratio and claim share are over property-years; the
rest over structures. Pre-FIRM structures were built before their community's first flood map;
FIRM status is the structure's last observed year. Build of 2026-10-15 (afternoon, with the permit status history).
Property value and the retrofit permit are observed only for structures matched to an ATTOM
parcel, and the permit only in county-years with a permit feed. 2023 dollars.
\end{minipage}
\end{table}
